\begin{center}
    \large
    \begin{singlespace}
        \textbf{\chineseTitle{}} \\[0.5cm]
    \end{singlespace}

    \begin{singlespace}
    學生：\studentCnName\hspace{2.5cm} 指導教授：\advisorCnName\hspace{0.1cm} 博士 \\
    \end{singlespace}

    \vspace{0.5cm}

    \begin{singlespace}
        國立陽明交通大學\\
        \NameofDepartmentInstituteCN\\[0.2cm]
    \end{singlespace}

    \textbf{摘\hspace{1cm}要} \\[0.5cm]
\end{center}

\normalsize

近似最近鄰搜尋（approximate nearest neighbor search）已成為向量資料庫與檢索、推薦系統中的核心元件。
當資料規模成長至億級，索引無法完整載入記憶體，磁碟式圖索引因而成為主流方案，但其查詢延遲往往受制於大量的隨機磁碟讀取。
現有做法可概分為兩類：以 DiskANN 為代表的圖式索引能在有限記憶體下維持搜尋品質，卻在圖走訪過程中付出可觀的輸入輸出成本；
以 Starling 為代表的頁面布局（page layout）最佳化則提高每次讀取的資訊密度，但其效益取決於靜態布局與預先配置的記憶體預算，
且在高並發的服務情境下容易出現長尾延遲。
我們觀察到，磁碟式搜尋的浪費同時存在於空間與時間兩個面向：圖走訪的熱點高度集中於少數節點，
而全域固定的搜尋長度（search list size）使得簡單查詢在結果早已收斂之後仍持續展開。
為此，我們提出 PaceANN，一個無須重建索引、可直接套用於既有 DiskANN 索引的磁碟式搜尋加速方案：
其一是分層快取（layered caching），第一層為以 medoid 為根的靜態廣度優先快取，
第二層為分片式 CLOCK 快取（sharded CLOCK cache），能在查詢期間動態累積熱點節點的鄰居列表，兩層皆未命中時則退回原始的磁碟讀取路徑，確保搜尋語義不變；
其二是自適應提前終止（adaptive early termination），在束搜尋（beam search）迴圈中依據前 $k$ 個結果的收斂狀態，
為每個查詢個別決定停止時機，使簡單查詢得以提早結束，困難查詢則繼續展開。
實驗結果顯示，在 SIFT100M 的伺服器模式（server mode）與相同記憶體預算下，PaceANN 於召回率約 95\%、並發度為 8 的工作點上，
相較於 Starling 可達 6.56 倍的查詢吞吐量、6.8 倍的 P99 延遲改善與 13.8 倍的 P999 延遲改善。

\vspace{1cm}

\textbf{關鍵字：}近似最近鄰搜尋, 磁碟式圖索引, 分層快取, 提前終止, 尾延遲
